\begin{titlepage}
\centering
\vspace*{2.5cm}
{\sffamily\bfseries\fontsize{40}{46}\selectfont\color{nomadbraun} NOMAD\par}
\vspace{0.4cm}
{\sffamily\bfseries\fontsize{24}{28}\selectfont Handbuch\par}
\vspace{0.8cm}
{\Large Ein Wissensserver für Zuhause, der auch ohne Internet funktioniert\par}
\vspace{2.5cm}
\begin{tikzpicture}
\node[draw=nomadbraun,line width=1.6pt,rounded corners=4pt,inner sep=14pt,fill=nomadhell,text width=11cm,align=flush left]{%
\textbf{Dieses Heft gehört zu:}\par\vspace{6pt}
Rechner / USB-SSD: \zeile[6.2cm]\par\vspace{8pt}
Aufbewahrungsort: \zeile[6.6cm]\par\vspace{8pt}
Eingerichtet am: \zeile[3cm]\hfill von: \zeile[3.2cm]};
\end{tikzpicture}
\vfill
{\small Stand Oktober 2026, Software-Version 1.35.2.\par
Inoffizielle deutsche Fassung von Project NOMAD (Crosstalk Solutions), Apache-Lizenz 2.0.\par
Quellen dieses Heftes: \texttt{github.com/huppiflupp/project-nomad-de}, Ordner \texttt{de/handbuch}.\par}
\end{titlepage}

\chapter*{So benutzen Sie dieses Heft}
\addcontentsline{toc}{chapter}{So benutzen Sie dieses Heft}
Dieses Heft ist zum \textbf{Ausdrucken} gedacht. Legen Sie es zur USB-SSD oder zum Notfall-Laptop. Wenn es darauf ankommt, hat man meistens alles vergessen, deshalb steht hier jeder Schritt, auch der kleinste.

\textbf{Was Sie brauchen:} einen Rechner (am besten mit Windows, den Sie nicht verändern müssen), eine USB-Festplatte oder USB-SSD mit mindestens 512~GB, eine Internetverbindung \emph{zum Einrichten} und etwas Zeit. Danach läuft alles ohne Internet.

\textbf{Wie das Heft aufgebaut ist:}
\begin{itemize}[leftmargin=1.2em,itemsep=2pt]
\item Kapitel~1 erklärt, was NOMAD ist und \emph{ab wann} es ohne Internet funktioniert.
\item Kapitel~2 bis~6 führen Sie von der leeren USB-SSD bis zum bestandenen Offline-Test.
\item Kapitel~7 ist die \textbf{Notfall-Karte}, eine Seite zum Heraustrennen.
\item Kapitel~8 rechnet den Stromverbrauch durch (Akku, Solar, Kosten).
\item Kapitel~9 enthält Tabellen zum Ausfüllen von Hand: Softwarestand, Inhalte, Zugangsdaten.
\item Kapitel~10 beschreibt die Notfall-KI, einen Sprachassistenten ohne Internet.
\end{itemize}

\kasten{nomadbraun}{Zeichen in diesem Heft}{%
\abhaken Kästchen zum Abhaken.\par
\taste{Enter} Eine Taste auf der Tastatur.\par
\menue{Einstellungen} Ein Knopf oder Menüpunkt auf dem Bildschirm.\par
\befehl{so sieht ein Befehl aus}
Befehle tippen Sie genau so ab, mit Leerzeichen, und bestätigen mit \taste{Enter}. Groß- und Kleinschreibung zählt.}

\hinweis{Ein Teil der Inhalte (Wikipedia, medizinische Nachschlagewerke, Karten) liegt \textbf{auf Englisch} vor. Die Oberfläche von NOMAD und dieses Heft sind deutsch. Mehr dazu im Anhang „Was NOMAD nicht kann“.}
